\documentclass[12pt]{article}

% Idioma y codificación
\usepackage[spanish, es-tabla]{babel}       %es-tabla para que se titule "Tabla"
\usepackage[utf8]{inputenc}

% Márgenes
\usepackage[a4paper,top=3cm,bottom=2.5cm,left=3cm,right=3cm]{geometry}

% Comentarios de bloque
\usepackage{verbatim}

% Paquetes de links
\usepackage[hidelinks]{hyperref}    % Permite enlaces
\usepackage{url}                    % redirecciona a la web

% Más opciones para enumeraciones
\usepackage{enumitem}

% Personalizar la portada
\usepackage{titling}


% Paquetes de tablas
\usepackage{multirow}


%------------------------------------------------------------------------

%Paquetes de figuras
\usepackage{caption}
\usepackage{subcaption} % Figuras al lado de otras
\usepackage{float}      % Poner figuras en el sitio indicado H.


% Paquetes de imágenes
\usepackage{graphicx}       % Paquete para añadir imágenes
\usepackage{transparent}    % Para manejar la opacidad de las figuras

% Paquete para usar colores
\usepackage[dvipsnames]{xcolor}
\usepackage{pagecolor}      % Para cambiar el color de la página

% Habilita tamaños de fuente mayores
\usepackage{fix-cm}


%------------------------------------------------------------------------

% Paquetes de matemáticas
\usepackage{mathtools, amsfonts, amssymb, mathrsfs}
\usepackage[makeroom]{cancel}     % Simplificar tachando
\usepackage{polynom}    % Divisiones y Ruffini
\usepackage{units} % Para poner fracciones diagonales con \nicefrac

\usepackage{pgfplots}   %Representar funciones
\pgfplotsset{compat=1.18}  % Versión 1.18

\usepackage{tikz-cd}    % Para usar diagramas de composiciones
\usetikzlibrary{calc}   % Para usar cálculo de coordenadas en tikz

%Definición de teoremas, etc.
\usepackage{amsthm}
%\swapnumbers   % Intercambia la posición del texto y de la numeración

\theoremstyle{plain}

\makeatletter
\@ifclassloaded{article}{
  \newtheorem{teo}{Teorema}[section]
}{
  \newtheorem{teo}{Teorema}[chapter]  % Se resetea en cada chapter
}
\makeatother

\newtheorem{coro}{Corolario}[teo]           % Se resetea en cada teorema
\newtheorem{prop}[teo]{Proposición}         % Usa el mismo contador que teorema
\newtheorem{lema}[teo]{Lema}                % Usa el mismo contador que teorema

\theoremstyle{remark}
\newtheorem*{observacion}{Observación}

\theoremstyle{definition}

\makeatletter
\@ifclassloaded{article}{
  \newtheorem{definicion}{Definición} [section]     % Se resetea en cada chapter
}{
  \newtheorem{definicion}{Definición} [chapter]     % Se resetea en cada chapter
}
\makeatother

\newtheorem*{notacion}{Notación}
\newtheorem*{ejemplo}{Ejemplo}
\newtheorem*{ejercicio*}{Ejercicio}             % No numerado
\newtheorem{ejercicio}{Ejercicio} [section]     % Se resetea en cada section


% Modificar el formato de la numeración del teorema "ejercicio"
\renewcommand{\theejercicio}{%
  \ifnum\value{section}=0 % Si no se ha iniciado ninguna sección
    \arabic{ejercicio}% Solo mostrar el número de ejercicio
  \else
    \thesection.\arabic{ejercicio}% Mostrar número de sección y número de ejercicio
  \fi
}


% \renewcommand\qedsymbol{$\blacksquare$}         % Cambiar símbolo QED
%------------------------------------------------------------------------

% Paquetes para encabezados
\usepackage{fancyhdr}
\pagestyle{fancy}
\fancyhf{}

\newcommand{\helv}{ % Modificación tamaño de letra
\fontfamily{}\fontsize{12}{12}\selectfont}
\setlength{\headheight}{15pt} % Amplía el tamaño del índice


%\usepackage{lastpage}   % Referenciar última pag   \pageref{LastPage}
\fancyfoot[C]{\thepage}

%------------------------------------------------------------------------

% Conseguir que no ponga "Capítulo 1". Sino solo "1."
\makeatletter
\@ifclassloaded{book}{
  \renewcommand{\chaptermark}[1]{\markboth{\thechapter.\ #1}{}} % En el encabezado
    
  \renewcommand{\@makechapterhead}[1]{%
  \vspace*{50\p@}%
  {\parindent \z@ \raggedright \normalfont
    \ifnum \c@secnumdepth >\m@ne
      \huge\bfseries \thechapter.\hspace{1em}\ignorespaces
    \fi
    \interlinepenalty\@M
    \Huge \bfseries #1\par\nobreak
    \vskip 40\p@
  }}
}
\makeatother

%------------------------------------------------------------------------
% Paquetes de cógido
\usepackage{minted}
\renewcommand\listingscaption{Código fuente}

\usepackage{fancyvrb}
% Personaliza el tamaño de los números de línea
\renewcommand{\theFancyVerbLine}{\small\arabic{FancyVerbLine}}

% Estilo para C++
\newminted{cpp}{
    frame=lines,
    framesep=2mm,
    baselinestretch=1.2,
    linenos,
    escapeinside=||
}



\usepackage{listings} % Para incluir código desde un archivo

\renewcommand\lstlistingname{Código Fuente}
\renewcommand\lstlistlistingname{Índice de Códigos Fuente}

% Definir colores
\definecolor{vscodepurple}{rgb}{0.5,0,0.5}
\definecolor{vscodeblue}{rgb}{0,0,0.8}
\definecolor{vscodegreen}{rgb}{0,0.5,0}
\definecolor{vscodegray}{rgb}{0.5,0.5,0.5}
\definecolor{vscodebackground}{rgb}{0.97,0.97,0.97}
\definecolor{vscodelightgray}{rgb}{0.9,0.9,0.9}

% Configuración para el estilo de C similar a VSCode
\lstdefinestyle{vscode_C}{
  backgroundcolor=\color{vscodebackground},
  commentstyle=\color{vscodegreen},
  keywordstyle=\color{vscodeblue},
  numberstyle=\tiny\color{vscodegray},
  stringstyle=\color{vscodepurple},
  basicstyle=\scriptsize\ttfamily,
  breakatwhitespace=false,
  breaklines=true,
  captionpos=b,
  keepspaces=true,
  numbers=left,
  numbersep=5pt,
  showspaces=false,
  showstringspaces=false,
  showtabs=false,
  tabsize=2,
  frame=tb,
  framerule=0pt,
  aboveskip=10pt,
  belowskip=10pt,
  xleftmargin=10pt,
  xrightmargin=10pt,
  framexleftmargin=10pt,
  framexrightmargin=10pt,
  framesep=0pt,
  rulecolor=\color{vscodelightgray},
  backgroundcolor=\color{vscodebackground},
}

%------------------------------------------------------------------------

% Comandos definidos
\newcommand{\bb}[1]{\mathbb{#1}}
\newcommand{\cc}[1]{\mathcal{#1}}

% I prefer the slanted \leq
\let\oldleq\leq % save them in case they're every wanted
\let\oldgeq\geq
\renewcommand{\leq}{\leqslant}
\renewcommand{\geq}{\geqslant}

% Si y solo si
\newcommand{\sii}{\iff}

% Letras griegas
\newcommand{\eps}{\epsilon}
\newcommand{\veps}{\varepsilon}
\newcommand{\lm}{\lambda}

\newcommand{\ol}{\overline}
\newcommand{\ul}{\underline}
\newcommand{\wt}{\widetilde}
\newcommand{\wh}{\widehat}

\let\oldvec\vec
\renewcommand{\vec}{\overrightarrow}

% Derivadas parciales
\newcommand{\del}[2]{\frac{\partial #1}{\partial #2}}
\newcommand{\Del}[3]{\frac{\partial^{#1} #2}{\partial^{#1} #3}}
\newcommand{\deld}[2]{\dfrac{\partial #1}{\partial #2}}
\newcommand{\Deld}[3]{\dfrac{\partial^{#1} #2}{\partial^{#1} #3}}


\newcommand{\AstIg}{\stackrel{(\ast)}{=}}
\newcommand{\Hop}{\stackrel{L'H\hat{o}pital}{=}}

\newcommand{\red}[1]{{\color{red}#1}} % Para integrales, destacar los cambios.

% Método de integración
\newcommand{\MetInt}[2]{
    \left[\begin{array}{c}
        #1 \\ #2
    \end{array}\right]
}

% Declarar aplicaciones
% 1. Nombre aplicación
% 2. Dominio
% 3. Codominio
% 4. Variable
% 5. Imagen de la variable
\newcommand{\Func}[5]{
    \begin{equation*}
        \begin{array}{rrll}
            #1:& #2 & \longrightarrow & #3\\
               & #4 & \longmapsto & #5
        \end{array}
    \end{equation*}
}

%------------------------------------------------------------------------
\let\oldRe\Re % save them in case they're every wanted
\let\oldIm\Im
\renewcommand{\Re}{\operatorname{Re}} % redefine them
\renewcommand{\Im}{\operatorname{Im}}
\DeclareMathOperator{\Log}{Log}
\DeclareMathOperator{\Arg}{Arg}

\begin{document}

    % 1. Foto de fondo
    % 2. Título
    % 3. Encabezado Izquierdo
    % 4. Color de fondo
    % 5. Coord x del titulo
    % 6. Coord y del titulo
    % 7. Fecha

    
    % 1. Foto de fondo
% 2. Título
% 3. Encabezado Izquierdo
% 4. Color de fondo
% 5. Coord x del titulo
% 6. Coord y del titulo
% 7. Fecha

\newcommand{\portada}[7]{

    \portadaBase{#1}{#2}{#3}{#4}{#5}{#6}{#7}
    \portadaBook{#1}{#2}{#3}{#4}{#5}{#6}{#7}
}

\newcommand{\portadaExamen}[7]{

    \portadaBase{#1}{#2}{#3}{#4}{#5}{#6}{#7}
    \portadaArticle{#1}{#2}{#3}{#4}{#5}{#6}{#7}
}




\newcommand{\portadaBase}[7]{

    % Tiene la portada principal y la licencia Creative Commons
    
    % 1. Foto de fondo
    % 2. Título
    % 3. Encabezado Izquierdo
    % 4. Color de fondo
    % 5. Coord x del titulo
    % 6. Coord y del titulo
    % 7. Fecha
    
    
    \thispagestyle{empty}               % Sin encabezado ni pie de página
    \newgeometry{margin=0cm}        % Márgenes nulos para la primera página
    
    
    % Encabezado
    \fancyhead[L]{\helv #3}
    \fancyhead[R]{\helv \nouppercase{\leftmark}}
    
    
    \pagecolor{#4}        % Color de fondo para la portada
    
    \begin{figure}[p]
        \centering
        \transparent{0.3}           % Opacidad del 30% para la imagen
        
        \includegraphics[width=\paperwidth, keepaspectratio]{assets/#1}
    
        \begin{tikzpicture}[remember picture, overlay]
            \node[anchor=north west, text=white, opacity=1, font=\fontsize{60}{90}\selectfont\bfseries\sffamily, align=left] at (#5, #6) {#2};
            
            \node[anchor=south east, text=white, opacity=1, font=\fontsize{12}{18}\selectfont\sffamily, align=right] at (9.7, 3) {\textbf{\href{https://losdeldgiim.github.io/}{Los Del DGIIM}}};
            
            \node[anchor=south east, text=white, opacity=1, font=\fontsize{12}{15}\selectfont\sffamily, align=right] at (9.7, 1.8) {Doble Grado en Ingeniería Informática y Matemáticas\\Universidad de Granada};
        \end{tikzpicture}
    \end{figure}
    
    
    \restoregeometry        % Restaurar márgenes normales para las páginas subsiguientes
    \pagecolor{white}       % Restaurar el color de página
    
    
    \newpage
    \thispagestyle{empty}               % Sin encabezado ni pie de página
    \begin{tikzpicture}[remember picture, overlay]
        \node[anchor=south west, inner sep=3cm] at (current page.south west) {
            \begin{minipage}{0.5\paperwidth}
                \href{https://creativecommons.org/licenses/by-nc-nd/4.0/}{
                    \includegraphics[height=2cm]{assets/Licencia.png}
                }\vspace{1cm}\\
                Esta obra está bajo una
                \href{https://creativecommons.org/licenses/by-nc-nd/4.0/}{
                    Licencia Creative Commons Atribución-NoComercial-SinDerivadas 4.0 Internacional (CC BY-NC-ND 4.0).
                }\\
    
                Eres libre de compartir y redistribuir el contenido de esta obra en cualquier medio o formato, siempre y cuando des el crédito adecuado a los autores originales y no persigas fines comerciales. 
            \end{minipage}
        };
    \end{tikzpicture}
    
    
    
    % 1. Foto de fondo
    % 2. Título
    % 3. Encabezado Izquierdo
    % 4. Color de fondo
    % 5. Coord x del titulo
    % 6. Coord y del titulo
    % 7. Fecha


}


\newcommand{\portadaBook}[7]{

    % 1. Foto de fondo
    % 2. Título
    % 3. Encabezado Izquierdo
    % 4. Color de fondo
    % 5. Coord x del titulo
    % 6. Coord y del titulo
    % 7. Fecha

    % Personaliza el formato del título
    \pretitle{\begin{center}\bfseries\fontsize{42}{56}\selectfont}
    \posttitle{\par\end{center}\vspace{2em}}
    
    % Personaliza el formato del autor
    \preauthor{\begin{center}\Large}
    \postauthor{\par\end{center}\vfill}
    
    % Personaliza el formato de la fecha
    \predate{\begin{center}\huge}
    \postdate{\par\end{center}\vspace{2em}}
    
    \title{#2}
    \author{\href{https://losdeldgiim.github.io/}{Los Del DGIIM}}
    \date{Granada, #7}
    \maketitle
    
    \tableofcontents
}




\newcommand{\portadaArticle}[7]{

    % 1. Foto de fondo
    % 2. Título
    % 3. Encabezado Izquierdo
    % 4. Color de fondo
    % 5. Coord x del titulo
    % 6. Coord y del titulo
    % 7. Fecha

    % Personaliza el formato del título
    \pretitle{\begin{center}\bfseries\fontsize{42}{56}\selectfont}
    \posttitle{\par\end{center}\vspace{2em}}
    
    % Personaliza el formato del autor
    \preauthor{\begin{center}\Large}
    \postauthor{\par\end{center}\vspace{3em}}
    
    % Personaliza el formato de la fecha
    \predate{\begin{center}\huge}
    \postdate{\par\end{center}\vspace{5em}}
    
    \title{#2}
    \author{\href{https://losdeldgiim.github.io/}{Los Del DGIIM}}
    \date{Granada, #7}
    \thispagestyle{empty}               % Sin encabezado ni pie de página
    \maketitle
    \vfill
}
    \portadaExamen{ffccA4.jpg}{Variable Compleja I\\Examen II}{Variable Compleja I. Examen II}{MidnightBlue}{-8}{28}{2024-2025}{Arturo Olivares Martos}

    \begin{description}
        \item[Asignatura] Variable Compleja I.
        \item[Curso Académico] 2023-24.
        \item[Grado] Doble Grado en Ingeniería Informática y Matemáticas.
        \item[Grupo] Único.
        \item[Profesor] Javier Merí de la Maza.
        \item[Descripción] Prueba Intermedia.
        \item[Fecha] 7 de Mayo de 2024.
        \item[Duración] 120 minutos.
    
    \end{description}
    \newpage

    \begin{ejercicio}[4 puntos]
        Probar que la serie $\displaystyle \sum_{n\geq 1} \frac{1}{n^z}$ converge absolutamente en todo punto del dominio $\Omega = \{ z \in \mathbb{C} \mid \Re z > 1 \}$ y uniformemente en cada subconjunto compacto contenido en $\Omega$. Deducir que siguiente función $g$ es continua en $\Omega$ y calcular $\displaystyle \int_{C(\pi,1)} g(z) \ dz$:
        \Func{g}{\Omega}{\mathbb{C}}{z}{\displaystyle \sum_{n=1}^{\infty} \frac{1}{n^z}}
    \end{ejercicio}

    \begin{ejercicio}[3 puntos]
        Estudiar la derivabilidad de la siguiente función $f$:
        \Func{f}{\mathbb{C}}{\mathbb{C}}{z}{z e^{\ol{z}}}
    \end{ejercicio}

    \begin{ejercicio}[3 puntos]~
        \begin{enumerate}
            \item \textbf{[1.5 puntos]} Calcular la siguiente integral:
            \begin{equation*}
                \int_{C(0,1)} \frac{\cos(z)}{z(z-3)^3} \, dz
            \end{equation*}
            \item \textbf{[1.5 puntos]} Sean $f$ y $g$ dos funciones enteras verificando $f(z) = g(z)$ para cada $z \in \bb{T}$. Demostrar que $f(z) = g(z)$ para cada $z \in \ol{D}(0, 1)$. 
            \item \textbf{[1.5 puntos Extra]} Probar que, de hecho, $f = g$.
        \end{enumerate}
    \end{ejercicio}


    \newpage
    \setcounter{ejercicio}{0}

    \begin{ejercicio}[4 puntos]
        Probar que la serie $\displaystyle \sum_{n\geq 1} \frac{1}{n^z}$ converge absolutamente en todo punto del dominio $\Omega = \{ z \in \mathbb{C} \mid \Re z > 1 \}$ y uniformemente en cada subconjunto compacto contenido en $\Omega$. Deducir que siguiente función $g$ es continua en $\Omega$ y calcular $\displaystyle \int_{C(\pi,1)} g(z) \ dz$:
        \Func{g}{\Omega}{\mathbb{C}}{z}{\displaystyle \sum_{n=1}^{\infty} \frac{1}{n^z}}

        Estudiamos en primer lugar la convergencia uniforme. Sea $K\subset \Omega$ compacto. Como la parte real de $z$ es continua, $\exists M\in \mathbb{R}$ tal que:
        \begin{equation*}
            M=\min\{\Re z : z\in K\} > 1
        \end{equation*}
        
        Entonces, para todo $z\in K$ se tiene que:
        \begin{align*}
            \left|\dfrac{1}{n^z}\right| &= \dfrac{1}{n^{\Re z}} \leq \dfrac{1}{n^M} \qquad \forall n\in \mathbb{N}
        \end{align*}

        Como $M> 1$, la serie $\displaystyle \sum_{n\geq 1} \frac{1}{n^M}$ converge, y por el Test de Weierstrass, la serie $\displaystyle \sum_{n\geq 1} \frac{1}{n^z}$ converge uniformemente en $K$.\\

        Para la convergencia absoluta, consideramos $z\in \Omega$ fijo. Como $\{z\}$ es compacto, por el Test de Weierstrass tenemos que converge absolutamente en $\{z\}$. Como $z$ es arbitrario, tenemos que la serie converge absolutamente en todo punto de $\Omega$.\\
        
        Probaremos ahora la continuidad de $g$ en $\Omega$, algo que no tenemos directo por no ser $\Omega$ compacto. Fijado $z\in \Omega$, por ser $\Omega$ abierto $\exists R\in \bb{R}^+$ tal que $D(z,R)\subset \Omega$. Entonces, tomando $r=\nicefrac{R}{2}$ se tiene que $z\in \ol{D}(z,r)\subset \Omega$, con $K=\ol{D}(z,r)$ compacto. Como $K$ es compacto y el término general de la serie de $g$ está formado por funciones continuas para todo $n\in \mathbb{N}$, la convergencia uniforme de la serie en $K$ nos garantiza que $g$ es continua en $z\in K$. Como $z$ era arbitrario en $\Omega$, se concluye que $g$ es continua en $\Omega$.\\

        Por último, calcularemos dicha integral. Como $\ol{D}(\pi,1)$ es compacto, la convergencia es uniforme. Por tanto, tenemos que:
        \begin{align*}
            \int_{C(\pi,1)} g(z) \ dz &= \int_{C(\pi,1)} \sum_{n=1}^{\infty} \frac{1}{n^z} \ dz = \sum_{n=1}^{\infty} \int_{C(\pi,1)} \frac{1}{n^z} \ dz \AstIg \sum_{n=1}^{\infty} 0 = 0
        \end{align*}
        donde en $(\ast)$ hemos usado que el integrando es holomorfo en $\Omega$, y como $\Omega$ es estrellado, por el Teorema Local de Cauchy, este admite primitiva en $\Omega$. Como $C(\pi,1)$ es un camino cerrado, dicha integral se anula.
    \end{ejercicio}

    \begin{ejercicio}[3 puntos]
        Estudiar la derivabilidad de la siguiente función $f$:
        \Func{f}{\mathbb{C}}{\mathbb{C}}{z}{z e^{\ol{z}}}

        \begin{description}
            \item[Ecuaciones de Cauchy-Riemann]~
            
            Con vistas a aplicar las ecuaciones de Cauchy-Riemann, definimos las siguientes funciones $u,v:\mathbb{R}^2\to\mathbb{R}$:
        \begin{align*}
            u(x,y) &= \Re(f(x+iy)) = \Re((x+iy)e^{x-iy}) = \Re((x+iy)(e^x(\cos y - i\sen y)))\\
            &= e^x\Re(x\cos y +y\sen y + i(y\cos y - x\sen y)) = e^x(x\cos y +y\sen y)\\
            v(x,y) &= \Im(f(x+iy)) = e^x(y\cos y - x\sen y)
        \end{align*}

        Calculamos ahora las derivadas parciales de $u$ y $v$:
        \begin{align*}
            \frac{\partial u}{\partial x}(x,y) &= e^x\left(x\cos y + y\sen y+\cos y\right)\\
            \frac{\partial u}{\partial y}(x,y) &= e^x\left(-x\sen y + \sen y + y\cos y\right)\\
            \frac{\partial v}{\partial x}(x,y) &= e^x\left(y\cos y -x\sen y -\sen y\right)\\
            \frac{\partial v}{\partial y}(x,y) &= e^x\left(-y\sen y + \cos y - x\cos y\right)
        \end{align*}

        La primera ecuación de Cauchy-Riemann nos dice:
        \begin{align*}
            \frac{\partial u}{\partial x}(x,y) &= \frac{\partial v}{\partial y}(x,y)\\
            \cancel{e^x}\left(x\cos y + y\sen y+\cancel{\cos y}\right) &= \cancel{e^x}\left(-y\sen y + \cancel{\cos y} - x\cos y\right)\\
            2y\sen y + 2x\cos y &= 0\\
            y\sen y + x\cos y &= 0
        \end{align*}

        La segunda ecuación de Cauchy-Riemann nos dice:
        \begin{align*}
            \frac{\partial u}{\partial y}(x,y) &= -\frac{\partial v}{\partial x}(x,y)\\
            \cancel{e^x}\left(-x\sen y + \cancel{\sen y} + y\cos y\right) &= -\cancel{e^x}\left(y\cos y - x\sen y - \cancel{\sen y}\right)\\
            2y\cos y - 2x\sen y &= 0\\
            y\cos y - x\sen y &= 0
        \end{align*}

        Como vemos, la resolución de dichas ecuaciones no es trivial, por lo que optamos por otro camino.


        \item[Otra forma]~
        
        Fijado $z\in \mathbb{C}$, veamos si $f$ es derivable en $z$. Para ello, distinguimos en función de si $z=0$ o no.
        \begin{itemize}
            \item \ul{Caso $z\neq 0$}:
            
            Sea $g$ la siguiente función:
            \Func{g}{\mathbb{C}}{\mathbb{C}}{z}{e^{\ol{z}}}

            Supuesto que $f$ sea derivable en $z$ entonces $g$ sería derivable en $z$, puesto que:
            \begin{equation*}
                g(z) = \frac{f(z)}{z}
            \end{equation*}

            Veamos no obstante que $g$ no es derivable en $z$.
            %\begin{description}
            %    \item[Ecuaciones de Cauchy-Riemann]~
                
                Definimos $u,v:\mathbb{R}^2\to\mathbb{R}$ como:
                \begin{align*}
                    u(x,y) &= \Re(g(x+iy)) = \Re(e^{x-iy}) = e^x\cos (-y)\\
                    v(x,y) &= \Im(g(x+iy)) = \Im(e^{x-iy}) = e^x\sen (-y)
                \end{align*}

                Calculamos ahora las derivadas parciales de $u$ y $v$:
                \begin{align*}
                    \frac{\partial u}{\partial x}(x,y) = e^x\cos (-y)&
                    \frac{\partial u}{\partial y}(x,y) = e^x\sen (-y)\\
                    \frac{\partial v}{\partial x}(x,y) = e^x\sen (-y)&
                    \frac{\partial v}{\partial y}(x,y) = -e^x\cos (-y)
                \end{align*}

                La primera ecuación de Cauchy-Riemann nos dice:
                \begin{align*}
                    \frac{\partial u}{\partial x}(x,y) = e^x\cos(-y)
                    = \frac{\partial v}{\partial y}(x,y) = -e^x\cos(-y)
                \end{align*}

                Para que esto se de, es necesario que $e^x=0$ (lo cual no se da) o que $\cos(-y)=\cos(y)=0$.

                La segunda ecuación de Cauchy-Riemann nos dice:
                \begin{align*}
                    \frac{\partial u}{\partial y}(x,y) = e^x\sen(-y)
                    = -\frac{\partial v}{\partial x}(x,y) = -e^x\sen(-y)
                \end{align*}
                Para que esto se de, es necesario que $e^x=0$ (lo cual no se da) o que $\sen(-y)=-\sen(y)=0$.

                Como no es posible que el seno y el coseno reales se anulen simultáneamente, se concluye que $g$ no es derivable en $z$.

                \begin{comment}
                \item[Otra forma]~
                
                Calculemos $\ol{g(z)}$:
                \begin{align*}
                    \ol{g(z)} &= \ol{e^{\ol{z}}} = e^{\Re z}\left(\cos(-\Im z) - i\sen(-\Im z)\right)
                    =\\&= e^{\Re z}\left(\cos(\Im z) + i\sen(\Im z)\right)
                    = e^z\qquad \forall z\in \mathbb{C}
                \end{align*}

                Por tanto, supuesto que $g$ es holomorfa en $\bb{C}^*$, como $\ol{g}$ también lo es, entonces 
                \end{comment}
            %\end{description}
            
            
            Por tanto, como $g$ no es derivable en $z$, $f$ tampoco lo es.

            \item{ \ul{Caso $z=0$}:}
            
            Lo calcularemos por la definición formal.
            \begin{equation*}
                \hspace{-2cm}f'(0) = \lim_{z\to 0} \frac{f(z)-f(0)}{z-0} = \lim_{z\to 0} \frac{z e^{\ol{z}}}{z} = \lim_{z\to 0} e^{\ol{z}} = \exp\left(\lim_{z\to 0} \ol{z}\right) = \exp\left(\ol{\lim_{z\to 0} z}\right) = \exp\left(\ol{0}\right) = \exp(0) = 1
            \end{equation*}
            donde hemos usado la continuidad de la exponencial y de la conjugación. Por tanto, $f$ es derivable en $0$ y $f'(0)=1$.
        \end{itemize}
        
        Por tanto, $f$ es derivable en $0$ y no es derivable en ningún otro punto de $\mathbb{C}$.        
        \end{description}
    \end{ejercicio}

    \begin{ejercicio}[3 puntos]~
        \begin{enumerate}
            \item \textbf{[1.5 puntos]} Calcular la siguiente integral:
            \begin{equation*}
                \int_{C(0,1)} \frac{\cos(z)}{z(z-3)^3} \, dz
            \end{equation*}

            Definimos la función $f$ como:
            \Func{f}{D(0,2)}{\mathbb{C}}{z}{\dfrac{\cos(z)}{(z-3)^3}}

            Como $f$ es racional, $f\in \cc{H}(D(0,2))$. Por tanto, por la fórmula de Cauchy para la circunferencia, tenemos que:
            \begin{align*}
                \int_{C(0,1)} \dfrac{f(z)}{z}\ dz &= 2\pi i \cdot f(0)=-\dfrac{2\pi}{27}\cdot i
            \end{align*}
            \item \textbf{[1.5 puntos]} Sean $f$ y $g$ dos funciones enteras verificando $f(z) = g(z)$ para cada $z \in \bb{T}$. Demostrar que $f(z) = g(z)$ para cada $z \in \ol{D}(0, 1)$.
            
            Como son funciones enteras, para cada $z\in D(0,1)$ se tiene por la fórmula de Cauchy para la circunferencia que:
            \begin{align*}
                f(z) &= \frac{1}{2\pi i} \int_{C(0,1)} \frac{f(w)}{w-z} \, dw
                = \frac{1}{2\pi i} \int_{C(0,1)} \frac{g(w)}{w-z} \, dw
                = g(z)
            \end{align*}
            Además, para cada $z\in \bb{T}$ también se tiene por hipótesis que $f(z)=g(z)$. Por tanto, $f(z)=g(z)$ para cada $z\in \ol{D}(0,1)$.
            \item \textbf{[1.5 puntos Extra]} Probar que, de hecho, $f = g$.
            
            Si consideramos las restricciones a $\Omega=\ol{D}(0,1)$, se tiene que:
            \begin{equation*}
                f_{\big| \Omega} = g_{\big| \Omega}
            \end{equation*}

            Por el carácter local de la derivabilidad, se tiene que:
            \begin{align*}
                f^{(n)}(0) = g^{(n)}(0) \qquad \forall n\in \mathbb{N}
            \end{align*}

            Por otro lado, como $f,g$ son funciones enteras, por el Teorema de Cauchy son analíticas en $\mathbb{C}$. De hecho, considerando el desarrollo de taylor centrado en el origen, se tiene que:
            \begin{align*}
                f(z) &= \sum_{n=0}^{\infty} \frac{f^{(n)}(0)}{n!} z^n = \sum_{n=0}^{\infty} \frac{g^{(n)}(0)}{n!} z^n = g(z)\qquad \forall z\in \bb{C}
            \end{align*}
        \end{enumerate}
    \end{ejercicio}
\end{document}